% Môn: Vật lý
% Lớp: 11
% Chương: Chương I. Dao động
% Bài: L11C1 Bài 2. Mô tả dao động điều hòa · Xác định các đại lượng đặc trưng của dao động dựa vào số dao động toàn phần và thời gian
% ===== Câu 27 =====
% ID: L11C1B2-10-TLN
% Mức: TH
\dangbt{Xác định các đại lượng đặc trưng của dao động dựa vào số dao động toàn phần và thời gian}
\begin{ex}
% ID: L11C1B2-10-TLN
% Mức: TH
Một vật thực hiện 12 dao động trong $6\,\text{s}$. Tần số dao động theo đơn vị Hz bằng bao nhiêu?
\shortans{2}
\end{ex}

% ===== Câu 31 =====
% ID: L11C1B2-11-TLN
% Mức: VD
\begin{ex}
% ID: L11C1B2-11-TLN
% Mức: VD
Một vật thực hiện 18 dao động trong $9\,\text{s}$. Tần số dao động theo đơn vị Hz bằng bao nhiêu?
\shortans{2}
\end{ex}

% ===== Câu 35 =====
% ID: L11C1B2-12-TLN
% Mức: VD
\begin{ex}
% ID: L11C1B2-12-TLN
% Mức: VD
Một vật thực hiện 24 dao động trong $12\,\text{s}$. Tần số dao động theo đơn vị Hz bằng bao nhiêu?
\shortans{2}
\end{ex}


% ===== Câu 72 =====
% ID: L11C1B2-16-DS
% Mức: VD
\begin{ex}
% ID: L11C1B2-16-DS
% Mức: VD
Một vật dao động điều hòa thực hiện được đúng $120$ dao động toàn phần trong khoảng thời gian $1$ phút. Khảo sát các đại lượng chu kì và tần số của dao động:
\choiceTF
{\True Tần số dao động của vật bằng $2$ Hz.}
{\True Chu kì dao động của vật bằng $0,5$ s.}
{\True Tần số góc của dao động đạt giá trị $4\pi$ rad/s.}
{Khoảng thời gian ngắn nhất giữa hai lần liên tiếp vật đi qua vị trí cân bằng là $0,5$ s.}
\loigiai{
\begin{itemchoice}
\item (Đúng) Tần số dao động của vật bằng $2$ Hz. (Vì): Tần số dao động của vật là $f = \dfrac{N}{\Delta t} = \dfrac{120}{60} = 2$ Hz
\item (Đúng) Chu kì dao động của vật bằng $0,5$ s. (Vì): Chu kì dao động của vật là $T = \dfrac{1}{f} = \dfrac{1}{2} = 0,5$ s
\item (Đúng) Tần số góc của dao động đạt giá trị $4\pi$ rad/s. (Vì): Tần số góc của dao động là $\omega = 2\pi f = 2\pi \cdot 2 = 4\pi$ rad/s
\item (Sai) Khoảng thời gian ngắn nhất giữa hai lần liên tiếp vật đi qua vị trí cân bằng là $0,5$ s. (Vì): Khoảng thời gian ngắn nhất giữa hai lần liên tiếp vật đi qua vị trí cân bằng là một nửa chu kì, $\dfrac{T}{2} = \dfrac{0,5}{2} = 0,25$ s
\end{itemchoice}
}
\end{ex}

% ===== Câu 84 =====
% ID: L11C1B2-112-DS
% Mức: VD
\begin{ex}
% ID: L11C1B2-112-DS
% Mức: VD
Một vật thực hiện dao động điều hòa trên quỹ đạo là một đoạn thẳng dài $10$ cm. Biết rằng trong khoảng thời gian $10$ giây vật thực hiện được đúng $20$ dao động toàn phần. Xét các thông số động học của vật:
\choiceTF
{Biên độ dao động của vật bằng $10$ cm.}
{\True Chu kì dao động của vật bằng $0,5$ s.}
{\True Quãng đường vật đi được trong một chu kì dao động đầy đủ bằng $20$ cm.}
{Tốc độ trung bình của vật trong một dao động toàn phần có giá trị bằng $40$ cm/s.}
\loigiai{
\begin{itemchoice}
\item (Sai) Biên độ dao động của vật bằng $10$ cm. (Vì): Chiều dài quỹ đạo là $L = 10$ cm, nên $2A = 10$ cm, suy ra biên độ dao động của vật là $A = 5$ cm (không phải $10$ cm
\item (Đúng) Chu kì dao động của vật bằng $0,5$ s. (Vì): Chu kì dao động của vật là $T = \frac{\Delta t}{N} = \frac{10 \text{ s}}{20} = 0,5$ s
\item (Đúng) Quãng đường vật đi được trong một chu kì dao động đầy đủ bằng $20$ cm. (Vì): Quãng đường vật đi được trong một chu kì dao động đầy đủ là $S = 4A = 4 \cdot 5 = 20$ cm
\item (Sai) Tốc độ trung bình của vật trong một dao động toàn phần có giá trị bằng $40$ cm/s. (Vì): (Sai) Tốc độ trung bình của vật trong một dao động toàn phần có giá trị bằng $40$ cm/s
\end{itemchoice}
}
\end{ex}

% ===== Câu 96 =====
% ID: L11C1B2-115-DS
% Mức: VD
\begin{ex}
% ID: L11C1B2-115-DS
% Mức: VD
Một chất điểm chuyển động tròn đều với tốc độ góc bằng $5\pi$ rad/s trên đường tròn có bán kính bằng $6$ cm. Hình chiếu của chất điểm này lên một đường kính nằm ngang dao động điều hòa. Khảo sát các thông số của hình chiếu:
\choiceTF
{\True Biên độ dao động của hình chiếu bằng $6$ cm.}
{\True Chu kì dao động của hình chiếu có giá trị bằng $0,4$ s.}
{\True Tần số dao động của hình chiếu bằng $2,5$ Hz.}
{Trong khoảng thời gian $1$ phút, hình chiếu thực hiện được đúng $100$ dao động toàn phần.}
\loigiai{
\begin{itemchoice}
\item (Đúng) Biên độ dao động của hình chiếu bằng $6$ cm. (Vì): Biên độ dao động của hình chiếu bằng bán kính quỹ đạo chuyển động tròn đều: $A = R = 6$ cm
\item (Đúng) Chu kì dao động của hình chiếu có giá trị bằng $0,4$ s. (Vì): Chu kì dao động của hình chiếu bằng chu kì quay của chất điểm: $T = \dfrac{2\pi}{\omega} = \dfrac{2\pi}{5\pi} = 0,4$ s
\item (Đúng) Tần số dao động của hình chiếu bằng $2,5$ Hz. (Vì): Tần số dao động của hình chiếu: $f = \dfrac{1}{T} = \dfrac{1}{0,4} = 2,5$ Hz
\item (Sai) Trong khoảng thời gian $1$ phút, hình chiếu thực hiện được đúng $100$ dao động toàn phần. (Vì): (Sai) Trong khoảng thời gian $1$ phút, hình chiếu thực hiện được đúng $100$ dao động toàn phần
\end{itemchoice}
}
\end{ex}

% ===== Câu 105 =====
% ID: L11C1B2-118-DS
% Mức: VD
\begin{ex}
% ID: L11C1B2-118-DS
% Mức: VD
Một vật dao động điều hòa có chu kì dao động xác định là $T = 2$ s. Đánh giá tính đúng/sai của các nhận định dưới đây:
\choiceTF
{\True Trong thời gian 1 phút, vật thực hiện được đúng 30 dao động toàn phần.}
{Thời gian ngắn nhất để vật đi từ vị trí biên âm sang vị trí biên dương là $2\,\text{s}$.}
{Tần số dao động của vật đạt giá trị bằng $2\,\text{Hz}$.}
{\True Tần số góc tương ứng của dao động là $\pi$ rad/s.}
\loigiai{
\begin{itemchoice}
\item (Đúng) Trong thời gian 1 phút, vật thực hiện được đúng 30 dao động toàn phần. (Vì): Số dao động toàn phần trong thời gian $t = 1$ phút $= 60$ s là $n = \dfrac{t}{T} = \dfrac{60}{2} = 30$ dao động
\item (Sai) Thời gian ngắn nhất để vật đi từ vị trí biên âm sang vị trí biên dương là $2\,\text{s}$. (Vì): Thời gian ngắn nhất để vật đi từ vị trí biên âm sang vị trí biên dương là nửa chu kì, tức $\dfrac{T}{2} = \dfrac{2}{2} = 1$ s
\item (Sai) Tần số dao động của vật đạt giá trị bằng $2\,\text{Hz}$. (Vì): Tần số dao động của vật là $f = \dfrac{1}{T} = \dfrac{1}{2} = 0,5\,\text{Hz}$
\item (Đúng) Tần số góc tương ứng của dao động là $\pi$ rad/s. (Vì): Tần số góc tương ứng của dao động là $\omega = \dfrac{2\pi}{T} = \dfrac{2\pi}{2} = \pi\,\text{rad/s}$
\end{itemchoice}
}
\end{ex}

% ===== Câu 110 =====
% ID: L11C1B2-121-TN
% Mức: NB
\begin{ex}
% ID: L11C1B2-121-TN
% Mức: NB
Trong 2 giây, một vật dao động điều hòa thực hiện 8 dao động toàn phần. Chu kì dao động là bao nhiêu?
\choice
{$T = 0{,}25$ s}
{\True $T = 0{,}25$ s}
{$T = 0{,}5$ s}
{$T = 1$ s}
\loigiai{
Thời gian cho 8 dao động toàn phần là 2 giây, nên chu kì:
$T = \frac{\text{thời gian}}{\text{số dao động toàn phần}} = \frac{2}{8} = 0,25\,\text{s}$.
Vậy chu kì là $0,25\,\text{s}$.
}
\end{ex}

% ===== Câu 116 =====
% ID: L11C1B2-124-TN
% Mức: NB
\begin{ex}
% ID: L11C1B2-124-TN
% Mức: NB
Một vật dao động điều hòa thực hiện 50 dao động toàn phần trong 10 giây. Tần số dao động là bao nhiêu?
\choice
{$f = 5$ Hz}
{\True $f = 5$ Hz}
{$f = 10$ Hz}
{$f = 2$ Hz}
\loigiai{
Tần số là số dao động hoàn thành trong một giây:
  \begin{aligned}
 f &= \frac{\text{số dao động toàn phần}}{\text{thời gian}} 
&= $\frac{50}{10}&= 5\\text{Hz}$.
 \end{aligned} 
 Vậy tần số dao động là $5\,\text{Hz}$.
}
\end{ex}

% ===== Câu 122 =====
% ID: L11C1B2-127-TN
% Mức: NB
\begin{ex}
% ID: L11C1B2-127-TN
% Mức: NB
Một con lắc dao động điều hòa, thực hiện 15 dao động toàn phần trong 6 giây. Tính chu kì dao động.
\choice
{$T = 0{,}2$ s}
{$T = 0{,}3$ s}
{\True $T = 0{,}4$ s}
{$T = 0{,}5$ s}
\loigiai{
Chu kì dao động:
$T = \frac{\text{thời gian}}{\text{số dao động toàn phần}} = \frac{6}{15} = 0,4\,\mathrm{s}$.
Vậy chu kì là $0,4\,\mathrm{s}$.
}
\end{ex}

% ===== Câu 128 =====
% ID: L11C1B2-130-TN
% Mức: VD
\begin{ex}
% ID: L11C1B2-130-TN
% Mức: VD
Một vật thực hiện 10 dao động toàn phần trong thời gian $4\,\text{s}$. Tần số và tần số góc của dao động là bao nhiêu?
\choice
{Tần số $2,5\,\text{Hz}$, tần số góc $5\pi$ rad/s}
{\True Tần số $2,5\,\text{Hz}$, tần số góc $5\pi$ rad/s}
{Tần số $10\,\text{Hz}$, tần số góc $20\pi$ rad/s}
{Tần số $0,4\,\text{Hz}$, tần số góc $0,8\pi$ rad/s}
\loigiai{
Tính tần số của dao động:
  \begin{aligned}
 f &= $\dfrac{n}{t}&=\dfrac{10}{4}&= 2,5\\text{Hz}$.
 \end{aligned} 
 Tính tần số góc:
  \begin{aligned}
 $\omega$ &= 2 $\pi$ f 
&= 2 $\pi\cdot2,5$ &= 5 $\pi\\text{rad/s}$.
 \end{aligned} 
 Chọn đáp án tương ứng.
}
\end{ex}

% ===== Câu 134 =====
% ID: L11C1B2-133-TN
% Mức: VD
\begin{ex}
% ID: L11C1B2-133-TN
% Mức: VD
Một vật dao động điều hòa với biên độ $10\,\text{cm}$. Trong 30 giây, vật thực hiện được 120 dao động toàn phần. Tần số góc của dao động là bao nhiêu?
\choice
{$\omega = 4\pi$ rad/s}
{$\omega = 6\pi$ rad/s}
{\True $\omega = 8\pi$ rad/s}
{$\omega = 10\pi$ rad/s}
\loigiai{
Tần số dao động:
  \begin{aligned}
 f &= $\frac{120}{30}&= 4\\text{Hz}$.
 \end{aligned} 
 Tần số góc:
  \begin{aligned}
 $\omega$ &= 2 $\pi$ f 
&= 2 $\pi\cdot$ 4 
&= 8 $\pi\\text{rad/s}$.
 \end{aligned} 
 Vậy tần số góc là $8\pi$ rad/s.
}
\end{ex}

% ===== Câu 140 =====
% ID: L11C1B2-136-TN
% Mức: VDC
\begin{ex}% ID: L11C1B2-136-TN
% Mức: VDC
Một vật dao động điều hòa. Trong 4 giây, vật thực hiện được 20 dao động toàn phần. Biết biên độ dao động là $3 \text{ cm}$. Tại thời điểm ban đầu ($t = 0$), vật đi qua vị trí cân bằng theo chiều dương. Viết phương trình li độ của vật.
\choice
{$x = 3\sin(10\pi t)$ cm}
{\True $x = 3\sin(10\pi t)$ cm}
{$x = 3\cos(10\pi t)$ cm}
{$x = 3\sin(5\pi t)$ cm}
\loigiai{
Đề bài cho biết trong thời gian $t = 4\text{ s}$, vật thực hiện được $N = 20$ dao động toàn phần.\n\nTa có tần số dao động $f = \frac{N}{t} = \frac{20}{4} = 5\text{ Hz}$.\n\nTần số góc của dao động là $\omega = 2\pi f = 2\pi \cdot 5 = 10\pi \text{ rad/s}$.\n\nBiên độ dao động là $A = 3\text{ cm}$.\n\nTại thời điểm ban đầu $t = 0$, vật đi qua vị trí cân bằng theo chiều dương. Điều này có nghĩa là li độ $x = 0$ và vận tốc $v > 0$.\n\nNếu chọn phương trình dao động dạng $x = A\cos(\omega t + \varphi)$, thì tại $t=0$: $x(0) = A\cos\varphi = 0 \Rightarrow \cos\varphi = 0 \Rightarrow \varphi = \pm \frac{\pi}{2}$.\n\nVận tốc $v = x' = -A\omega\sin(\omega t + \varphi)$. Tại $t=0$: $v(0) = -A\omega\sin\varphi > 0$. Vì $A > 0, \omega > 0$, nên $\sin\varphi < 0$. Từ đó suy ra $\varphi = \frac{\pi}{2}$ (loại) hoặc $\varphi = -\frac{\pi}{2}$.\n\nVậy phương trình dao động là $x = 3\cos(10\pi t - \frac{\pi}{2})\text{ cm}$.\n\nSử dụng công thức lượng giác $\cos(\alpha - \frac{\pi}{2}) = \sin\alpha$, ta có $x = 3\sin(10\pi t)\text{ cm}$.\n\nVậy phương trình li độ của vật là $x = 3\sin(10\pi t)\text{ cm}$.
}
\end{ex}
